\documentclass[12pt]{article}
\usepackage[utf8]{inputenc}
\usepackage[T1]{fontenc}
\usepackage{hyperref}
\title{Yangbye Language Resources for K-12 Curriculum Development}
\author{Compiled by Claude Sonnet 5.5 (claude-sonnet-5-5) for Andr\'as Kornai}
\date{}
\begin{document}
\maketitle

\section{Introduction}
Yangbye (ISO 639-3 ybd; Glottocode yang1301) is a retired code. The ISO registration merged it into Rakhine (rki) in 2012. Yangbye is also called Ramree and is a dialect of Rakhine, spoken in southern Rakhine State, Myanmar. Wordlists show at least 94\% lexical similarity with the prestige Rakhine of Sittwe. The Glottolog family label ``Bookkeeping'' reflects its status as a retired code. Resources for Yangbye should therefore be sought under Rakhine. This search did not look for Rakhine materials.

\section{Primary Educational Resources}
None found for Yangbye.

\section{Linguistic and Reference Materials}
\begin{itemize}
\item Ramree dialect (Wikipedia). \url{https://en.wikipedia.org/wiki/Ramree_dialect}
\item IANA subtag record for ybd. \url{https://www.iana.org/assignments/lang-subtags-templates/ybd.txt}
\item Glottolog entry for Yangbye. \url{https://glottolog.org/resource/languoid/id/yang1301}
\end{itemize}

\section{Religious and Community Resources}
Scripture Earth and Joshua Project list Rakhine, not Yangbye. \url{https://www.scriptureearth.org/00eng.php?iso=rki} and \url{https://www.joshuaproject.net/index.php/languages/rki}

\section{Digital and Audio Resources}
None found for Yangbye.

\section{Curriculum Development Context}
No Yangbye-specific resources exist in the search results. A curriculum project should treat this as Rakhine (rki) and search under that code.

\end{document}